% Options for packages loaded elsewhere
\PassOptionsToPackage{unicode}{hyperref}
\PassOptionsToPackage{hyphens}{url}
%
\documentclass[
  11pt,
]{article}
\usepackage{amsmath,amssymb}
\usepackage{setspace}
\usepackage{iftex}
\ifPDFTeX
  \usepackage[T1]{fontenc}
  \usepackage[utf8]{inputenc}
  \usepackage{textcomp} % provide euro and other symbols
\else % if luatex or xetex
  \usepackage{unicode-math} % this also loads fontspec
  \defaultfontfeatures{Scale=MatchLowercase}
  \defaultfontfeatures[\rmfamily]{Ligatures=TeX,Scale=1}
\fi
\usepackage{lmodern}
\ifPDFTeX\else
  % xetex/luatex font selection
\fi
% Use upquote if available, for straight quotes in verbatim environments
\IfFileExists{upquote.sty}{\usepackage{upquote}}{}
\IfFileExists{microtype.sty}{% use microtype if available
  \usepackage[]{microtype}
  \UseMicrotypeSet[protrusion]{basicmath} % disable protrusion for tt fonts
}{}
\makeatletter
\@ifundefined{KOMAClassName}{% if non-KOMA class
  \IfFileExists{parskip.sty}{%
    \usepackage{parskip}
  }{% else
    \setlength{\parindent}{0pt}
    \setlength{\parskip}{6pt plus 2pt minus 1pt}}
}{% if KOMA class
  \KOMAoptions{parskip=half}}
\makeatother
\usepackage{xcolor}
\usepackage[margin=0.9in]{geometry}
\setlength{\emergencystretch}{3em} % prevent overfull lines
\providecommand{\tightlist}{%
  \setlength{\itemsep}{0pt}\setlength{\parskip}{0pt}}
\setcounter{secnumdepth}{-\maxdimen} % remove section numbering
\usepackage{microtype}
\usepackage{xurl}
\setlength{\emergencystretch}{3em}
\ifLuaTeX
  \usepackage{selnolig}  % disable illegal ligatures
\fi
\usepackage{bookmark}
\IfFileExists{xurl.sty}{\usepackage{xurl}}{} % add URL line breaks if available
\urlstyle{same}
\hypersetup{
  hidelinks,
  pdfcreator={LaTeX via pandoc}}

\author{}
\date{}

\begin{document}

\setstretch{1.03}
\section{Quadratic-form density: absolute continuity and a
separable-perturbation
theorem}\label{quadratic-form-density-absolute-continuity-and-a-separable-perturbation-theorem}

22 September 2026

\textbf{Scope.} This note proves absolute continuity for every nonzero
quadratic form under the stated product-density assumptions, recalls the
exact diagonal bound, and proves a bounded-continuous-density theorem
for a diffuse quadratic component with arbitrary separable
perturbations. It does not prove the full three-large-eigenvalue
conjecture for every symmetric matrix. The new extension below has a
complete analytic argument but has not yet received independent
end-to-end review. No novelty claim relative to the entire literature is
made.

\subsection{1. The full target}\label{the-full-target}

Let \(X_1,\ldots,X_n\) be independent real random variables with
probability densities \(f_i\) satisfying \(0\le f_i\le1\) almost
everywhere. For every real symmetric matrix A with at least three
eigenvalues of absolute value at least 1, the target is one universal
finite C such that

\[
\mathbb P\{X^TAX\in[t,t+h]\}\le Ch\qquad(t\in\mathbb R,\ h>0).
\]

Equivalently, the law has an essentially bounded density at most C.
Absolute continuity alone, a diagonal special case, and an estimate
assuming a special eigenvector geometry are strictly weaker conclusions.

\subsection{2. Absolute continuity for every nonzero symmetric
matrix}\label{absolute-continuity-for-every-nonzero-symmetric-matrix}

\textbf{Proposition 1.} Suppose the \(X_i\) are independent and
absolutely continuous. If A is real symmetric and nonzero, then X\^{}TAX
is absolutely continuous. No bound on the input densities,
three-eigenvalue assumption, or moment hypothesis is needed.

\textbf{Proof.} First suppose \(a_ii\) is nonzero for some i.
Conditional on all coordinates except \(X_i\), the quadratic form is

\[
a_{ii}X_i^2+bX_i+c
\]

with deterministic b,c.~A nonconstant one-variable quadratic maps every
absolutely continuous probability law to an absolutely continuous law.
Indeed, split its domain at its unique critical point; on each open
monotonicity interval the inverse is continuously differentiable
locally. Exhausting each interval by compact subintervals away from the
critical point shows that the preimage of any Lebesgue-null set is null.
The critical point itself has probability zero. Every conditional law is
therefore absolutely continuous, and their mixture is absolutely
continuous.

Now suppose every diagonal entry is zero. Choose i,j with \(a_ij\)
nonzero. Write

\[
X^TAX=L(X_{-i})X_i+R(X_{-i}),\qquad
L(X_{-i})=2\sum_{k\ne i}a_{ik}X_k.
\]

The random variable L is absolutely continuous: condition on all its
summands except the nonzero multiple 2a\_ij \(X_j\). In particular
P(L=0)=0. Almost every conditional law given X\_\{-i\} is consequently a
nonconstant affine image of the absolutely continuous law of \(X_i\).
Its mixture is absolutely continuous. This covers every nonzero
symmetric A. QED.

The coarea representation, when desired, is

\[
p_Q(t)=\int_{\{x:x^TAx=t,\ Ax\ne0\}}
\frac{\prod_i f_i(x_i)}{2\lVert Ax\rVert_2}\,d\mathcal H^{n-1}(x)
\quad\text{for a.e. }t.
\]

The critical set ker(A) is Lebesgue-null. This identity identifies the
remaining integral to be bounded; it supplies no universal bound by
itself.

\subsection{3. Diagonal forms: the bound
pi}\label{diagonal-forms-the-bound-pi}

\textbf{Proposition 2.} If A is diagonal and at least three diagonal
coefficients have magnitude at least 1, then its output density is
bounded by pi. The density has a continuous version.

\textbf{Proof.} For a\textgreater0 and a random variable X with density
f\textless=1, the variable Y=a(X-b)\^{}2 has density, for
y\textgreater0,

\[
g_{a,b}(y)=\frac{f(b+\sqrt{y/a})+f(b-\sqrt{y/a})}{2\sqrt{ay}}
\le\frac{1}{\sqrt{ay}}.
\]

Two independent positive square terms with coefficients a,b\textgreater0
therefore have convolution density at most

\[
\frac1{\sqrt{ab}}\int_0^s\frac{du}{\sqrt{u(s-u)}}
=\frac{\pi}{\sqrt{ab}}\quad(s>0).
\]

Among three coefficients of magnitude at least 1, two have the same
sign. Reflect the sum when both are negative. Their density is bounded
by pi. Convolution with the independent remaining terms cannot increase
this bound.

For continuity, first convolve the bounded density of the same-sign pair
with the \(L^1\) density of the third nonzero square term. If g is
bounded and h is integrable, then

\[
\sup_t |(g*h)(t+v)-(g*h)(t)|\le\lVert g\rVert_\infty\lVert h(\cdot+v)-h\rVert_1\to0.
\]

Thus this convolution has a uniformly continuous version. Its
integrability and nonnegativity imply that it vanishes at infinity.
Further convolution with a probability measure preserves boundedness,
continuity, and vanishing at infinity. QED.

This argument uses independence in the original coordinates.
Diagonalizing a general symmetric matrix does not preserve that
independence.

\subsection{4. A bounded-density theorem stable under all separable
perturbations}\label{a-bounded-density-theorem-stable-under-all-separable-perturbations}

Let U be an n by r real matrix with orthonormal columns, where
\(r\ge3\). Denote its rows by \(u_i\)\^{}T and put

\[
c_i=\lVert u_i\rVert_2^2,\qquad
\sum_i u_i u_i^T=I_r,\qquad\sum_i c_i=r.
\]

Assume \(c_i\le1/2\) for every i. This is an additional geometric
hypothesis: no input coordinate accounts for more than half the squared
projection in this r-dimensional component. Let \(\Lambda\) be any
invertible real symmetric r by r matrix and set

\[
B(U)=\prod_{i:c_i>0}(1-c_i)^{-(1-c_i)/2},\qquad
\Delta=|\det\Lambda|.
\]

\textbf{Theorem 3 (separable-perturbation extension).} For independent
\(X_i\) with densities \(0\le f_i\le1\) and arbitrary real-valued
measurable functions \(V_i\), the variable

\[
T=(U^TX)^T\Lambda(U^TX)+\sum_{i=1}^n V_i(X_i)
\]

has a bounded continuous density, vanishing at infinity, and

\[
\boxed{\quad
\lVert p_T\rVert_\infty
\le \frac{r}{r-2}\,B(U)^{2/r}\,\Delta^{-1/r}
\le \frac{er}{r-2}\,\Delta^{-1/r}.
\quad}
\]

The \(V_i\) need only be finite almost surely. They need not be smooth,
bounded, polynomial, or integrable. In particular they can be arbitrary
diagonal quadratic terms plus arbitrary linear terms. The theorem
imposes no sign restriction on \(\Lambda\). If \(r=3\) and
\(\Delta\ge1\), the bound is at most 3e.

\subsubsection{4.1. The two classical analytic
inputs}\label{the-two-classical-analytic-inputs}

Use the Fourier convention

\[
\widehat h(\xi)=\int_{\mathbb R}e^{-i\xi x}h(x)\,dx.
\]

The geometric rank-one Brascamp-Lieb inequality states that, for unit
vectors \(v_i\) and \(c_i\)\textgreater0 with
\(\sum_i c_i v_i v_i^T=I_r\),

\[
\int_{\mathbb R^r}\prod_i H_i(v_i\cdot\xi)^{c_i}\,d\xi
\le\prod_i\left(\int_{\mathbb R}H_i(s)\,ds\right)^{c_i}
\]

for nonnegative integrable \(H_i\). The precise statement also appears
as Theorem 3.3 of Livshyts-Paouris-Pivovarov {[}2{]}.

Beckner's sharp Hausdorff-Young inequality {[}1{]}, valid for
complex-valued h, states that for \(1\le p\le2\) and \(q=p/(p-1)\),

\[
\lVert\widehat h\rVert_q\le
(2\pi)^{1/q}p^{1/(2p)}q^{-1/(2q)}\lVert h\rVert_p.
\]

These are the only non-elementary inequalities used in Theorem 3. The
remaining steps, including the complex-amplitude issue and the
oscillatory interchange, are proved below.

\subsubsection{4.2. A phase-uniform Fourier
bound}\label{a-phase-uniform-fourier-bound}

For a fixed real t, define complex input functions

\[
h_{i,t}(x)=f_i(x)e^{itV_i(x)}.
\]

For every \(p\ge1\), their norms satisfy

\[
\lVert h_{i,t}\rVert_p=\lVert f_i\rVert_p\le1.
\]

Let \(nu_t\) be the pushforward under x -\textgreater{} U\^{}T x of the
complex measure with density \(product_i\) h\_\{i,t\}(\(x_i\)). Since
U\^{}T has full row rank, this measure has a complex \(L^1\) density
\(g_t\). To see this directly, complete U to an orthogonal matrix (U,W)
and integrate the original density at Uy+Wz over z. Fubini gives

\[
\lVert g_t\rVert_1\le1.
\]

Its Fourier transform is

\[
\widehat g_t(\xi)=\prod_i\widehat h_{i,t}(u_i\cdot\xi).
\]

Zero rows contribute scalar factors of modulus at most 1 and may be
omitted in estimates. For \(c_i\)\textgreater0 set
\(v_i\)=\(u_i\)/sqrt(\(c_i\)). Brascamp-Lieb, with

\[
H_i(s)=|\widehat h_{i,t}(\sqrt{c_i}s)|^{1/c_i},
\]

gives

\[
\lVert\widehat g_t\rVert_1
\le\prod_{i:c_i>0}c_i^{-c_i/2}
\lVert\widehat h_{i,t}\rVert_{1/c_i}.
\tag{4.1}
\]

Set \(q_i\)=1/\(c_i\) and \(p_i\)=1/(1-\(c_i\)). Because \(c_i\le1/2\),
\(q_i\)\textgreater=2 and \(p_i\)\textless=2, so sharp Hausdorff-Young
applies:

\[
\lVert\widehat h_{i,t}\rVert_{1/c_i}
\le(2\pi)^{c_i}(1-c_i)^{-(1-c_i)/2}c_i^{c_i/2}.
\]

The powers of \(c_i\) cancel in (4.1), leaving

\[
\boxed{\lVert\widehat g_t\rVert_1\le(2\pi)^rB(U).}
\tag{4.2}
\]

This bound is uniform in t and in every \(V_i\). Finally,

\[
-(1-c)\log(1-c)\le c\quad(0\le c<1)
\]

and \(\sum_i c_i=r\) imply

\[
B(U)\le e^{r/2}.
\tag{4.3}
\]

\subsubsection{4.3. Gaussian-regularized quadratic
integration}\label{gaussian-regularized-quadratic-integration}

The characteristic function of T is exactly

\[
\varphi_T(t)=\int_{\mathbb R^r}e^{it y^T\Lambda y}g_t(y)\,dy.
\]

There is no claim that the separable perturbation is independent of
U\^{}T X. Its entire dependence was retained in the complex density
\(g_t\).

For epsilon\textgreater0 and t nonzero, introduce

\[
I_\varepsilon(t)=\int e^{it y^T\Lambda y-\varepsilon|y|^2}g_t(y)\,dy.
\]

By (4.2), Fourier inversion for \(g_t\) and absolutely justified Fubini
yield

\[
I_\varepsilon(t)=(2\pi)^{-r}\int\widehat g_t(\xi)K_{\varepsilon,t}(\xi)\,d\xi,
\]

where

\[
K_{\varepsilon,t}(\xi)=\int e^{-\varepsilon|y|^2+it y^T\Lambda y+i\xi\cdot y}\,dy.
\]

Diagonalize \(\Lambda\) orthogonally inside this Gaussian integral. If
its eigenvalues are \(lambda_j\) and the rotated frequency is eta, the
j-th factor is

\[
\sqrt\pi\,(\varepsilon-it\lambda_j)^{-1/2}
\exp\!\left(-\frac{\eta_j^2}{4(\varepsilon-it\lambda_j)}\right).
\]

Its exponential has modulus at most 1 because

\[
\operatorname{Re}(\varepsilon-it\lambda_j)^{-1}
=\frac{\varepsilon}{\varepsilon^2+t^2\lambda_j^2}>0.
\]

Consequently,

\[
|K_{\varepsilon,t}(\xi)|
\le\pi^{r/2}\prod_j(\varepsilon^2+t^2\lambda_j^2)^{-1/4}
\le\pi^{r/2}|t|^{-r/2}\Delta^{-1/2}.
\]

Combining this with (4.2) gives

\[
|I_\varepsilon(t)|\le B(U)\pi^{r/2}|t|^{-r/2}\Delta^{-1/2}.
\]

For each fixed t, \(g_t\) is integrable, so dominated convergence as
epsilon decreases to zero yields

\[
|\varphi_T(t)|\le\min\{1,K|t|^{-r/2}\},
\qquad K=B(U)\pi^{r/2}\Delta^{-1/2}.
\tag{4.4}
\]

This is not a stationary-phase estimate for arbitrary one-dimensional
positive densities. It relies essentially on the diffuse-frame
hypothesis and a phase-uniform, frequency-integrated product estimate.

\subsubsection{4.4. Integration and density
inversion}\label{integration-and-density-inversion}

Let \(t_0\)=K\^{}\{2/r\}. Since r\textgreater2,

\[
\begin{aligned}
\int_{\mathbb R}|\varphi_T(t)|\,dt
&\le2\left[t_0+K\int_{t_0}^\infty t^{-r/2}\,dt\right]\\
&=2\frac{r}{r-2}K^{2/r}.
\end{aligned}
\]

Thus \(phi_T\) belongs to \(L^1\)(R). Fourier inversion gives a bounded
continuous density vanishing at infinity and

\[
\lVert p_T\rVert_\infty
\le\frac1{2\pi}\lVert\varphi_T\rVert_1
\le\frac{r}{r-2}B(U)^{2/r}\Delta^{-1/r}.
\]

Equation (4.3) supplies the simpler bound er/(r-2) times
\(\Delta^{-1/r}\). This proves Theorem 3. QED.

\subsection{5. Consequences for matrices with unrestricted remaining
terms}\label{consequences-for-matrices-with-unrestricted-remaining-terms}

\textbf{Corollary 4 (diagonal perturbations).} Theorem 3 applies to

\[
A=D+U\Lambda U^T
\]

for every real diagonal D, with the same constant, and it allows any
linear term in Q. The spectrum of the full matrix A need not have rank r
or be bounded away from zero. The determinant in the bound is that of
the specified component \(\Lambda\), not the product of eigenvalues of
A.

\textbf{Corollary 5 (a principal block suffices).} Let I be a set of
original coordinate indices. Suppose the physical principal block
\(A_II\) has a representation

\[
A_{II}=D+U\Lambda U^T
\]

satisfying Theorem 3. Every other block of A, including the couplings
between I and its complement, may be arbitrary. Then X\^{}TAX has the
same bounded-continuous-density estimate.

\textbf{Proof.} Conditional on X\_\{I\^{}c\}=z, the form is

\[
X_I^TA_{II}X_I+2X_I^TA_{I,I^c}z+z^TA_{I^c,I^c}z.
\]

The conditional \(X_I\) retain their original independent densities.
Absorb the displayed linear term and the diagonal quadratic terms into
the \(V_i\). Theorem 3 provides (4.4) uniformly in z; a constant phase
has modulus one. Averaging the conditional characteristic functions
preserves (4.4), so the same \(L^1\) inversion argument proves the
corollary. QED.

For input bounds \(f_i\le M_i\) instead of 1, the same calculation
replaces B(U) by B(U) \(product_i\) \(M_i\)\^{}(\(c_i\)). In particular
a common bound M introduces exactly the expected \(M^2\) factor in the
final density estimate.

\subsubsection{A concrete full-rank indefinite
example}\label{a-concrete-full-rank-indefinite-example}

Take \(n=6\), \(r=3\),

\[
U=\frac1{\sqrt2}\begin{pmatrix}I_3\\I_3\end{pmatrix},\qquad
\Lambda=\operatorname{diag}(2,2,-2),\qquad D=\tfrac1{10}I_6.
\]

Every leverage equals 1/2. The full matrix A=D+U \(\Lambda\) U\^{}T has
eigenvalues

\[
21/10,\quad21/10,\quad-19/10,\quad1/10,\quad1/10,\quad1/10.
\]

It meets the original three-large-eigenvalue condition, is not diagonal
in the independent input coordinates, is indefinite, and has three
additional small nonzero eigenvalues. Here B(U)=\(2^{3/2}\),
\(\Delta\)=8, so Theorem 3 gives

\[
\lVert p_{X^TAX}\rVert_\infty\le3.
\]

The example illustrates a covered family, not a universal matrix
reduction.

\subsection{6. What remains for the full
conjecture}\label{what-remains-for-the-full-conjecture}

The extension removes a real limitation of the earlier diffuse-rank
theorem: a diagonal or separable remainder is not independent noise, yet
it can be retained exactly without worsening the bound. The
principal-block corollary also permits arbitrary couplings outside a
suitable physical block.

Two gaps remain. First, the spectral hypothesis alone does not establish
a representation A=D+U \(\Lambda\) U\^{}T with all leverage scores at
most 1/2 and determinant bounded below. Second, no argument above covers
every leverage-heavy geometry or an arbitrary nonseparable remainder by
simply adding its phase. For \(c_i>1/2\) the attempted Fourier step
would require an \(L^q\) estimate with \(q<2\), outside the available
Hausdorff-Young range. A new inequality or a genuine product-preserving
reduction is needed there.

The full target therefore remains unproved in this note. No conditional
implication is being presented as an unconditional solution.

\subsection{7. Sources and provenance}\label{sources-and-provenance}

{[}1{]} W. Beckner, \emph{Inequalities in Fourier analysis}, Annals of
Mathematics 102 (1975), 159-182. DOI: 10.2307/1970980. Publisher record:
https://annals.math.princeton.edu/1975/102-1/p11 . The precise sharp
Hausdorff-Young inequality is the imported classical input.

{[}2{]} G. Livshyts, G. Paouris, P. Pivovarov, \emph{On sharp bounds for
marginal densities of product measures}, arXiv:1507.07949. Theorem 3.3
states the normalized geometric Brascamp-Lieb inequality used above.
https://arxiv.org/pdf/1507.07949 . Their marginal-density theorem alone
is not a proof of the nonlinear target.

{[}3{]} The existing diffuse-rank argument being extended: private
research archive \path{yspennstate/math-open-problems}.\par
Commit \path{0b5925a458d36c1ec1cf8ab344c9676260a97c8a}.\par
File \path{proofnet24/QFORM-density/fleet-2026-09-22/fleet02/B08/B08-20260922T113624Z-qform11/REPORT.md},
blob \path{212febc02c4c1f3a8e6c073d0a8ed5d2d6b4ea0a}. Its unperturbed
Fourier/Brascamp-Lieb calculation was read directly. The new step is the
complex product amplitude retaining arbitrary separable phases, followed
by the principal-block conditioning corollary.

{[}4{]} Exact target, diagonal control, and pointwise-decay correction:
unpublished coordination notes, 22 September 2026.\par

A finite algebraic sanity check accompanies this note; it verifies the
explicit example and scalar prefactors only. It is not a proof of the
universal conjecture or a substitute for review of Theorem 3. This
interactive derivation is not an independent scheduled-worker review and
is not native learned-memory ingestion.

\end{document}
